% ===========================================================================
%  方法一览（自总结版：把「要点」一栏留空，或在 main.tex 注释本文件）
% ===========================================================================
\section*{方法一览}

\begin{center}
\begin{tabularx}{\linewidth}{@{}l L l@{}}
\toprule
\multicolumn{1}{l}{\color{indigo}方法} &
\multicolumn{1}{l}{\color{indigo}要点} &
\multicolumn{1}{l}{\color{indigo}出处} \\
\midrule
作差法 & 作差、通分、判断符号（糖水不等式） & 0.1、题 1 \\
\starref{} $\to$ 基本不等式 & 平方非负 $\Rightarrow a^2+b^2\ge 2ab$ $\Rightarrow$ 用 $\sqrt a$ 替换 $a$ & 卡片 1 \\
几何解释 & 半弦不过半径；一张图呈现四个平均数 & 题 5、拓展 A \\
证明不等式 & 逐项放缩，同向相乘，取等条件须一致（会师） & 题 6、7 \\
求最值 & 一正、二定、三相等 & 题 8--13 \\
配凑 & 改写为「分母 + 分母的倒数 + 常数」；配系数凑定值 & 题 10、12，练习 9 \\
条件最值 & 乘「1」代换；和积互化；双换元 & 题 14--17 \\
齐次化 & 变形围绕次数进行；0 次齐次式只依赖比值 & 拓展 C \\
对称性 & 用于猜测取等点，不能代替证明 & 拓展 D \\
\bottomrule
\end{tabularx}
\end{center}

\begin{strand}
学生自总结版：把「要点」一栏留空，听完课自行填写。
\end{strand}
